\section*{Background and Motivation}
Unmanned Aerial Vehicles (UAVs) have become important operational platforms in mission-critical settings such as disaster response, search-and-rescue, infrastructure inspection, environmental monitoring, and surveillance. Their practical value lies in rapid deployment, wide-area aerial sensing, and the ability to reach locations that may be inaccessible or unsafe for ground teams. These characteristics make UAVs especially attractive for time-sensitive missions in which real-time situational awareness and adaptive decision-making are essential \cite{mcenroe2022survey,baktayan2024slr}.

Within this broad landscape, small UAVs are particularly relevant because they are less costly, easier to deploy, and more agile than larger aerial platforms, making them well suited to cooperative swarm operation. Their mobility and scalability support distributed sensing and coordinated coverage, while recent advances in onboard perception have enabled AI-driven tasks such as object detection, tracking, mapping, and persistent monitoring. As a result, small UAV swarms are increasingly viewed as a practical substrate for edge intelligence in dynamic field environments \cite{jung2024edgeai,qu2024elastic,clerigues2024mission}.

The central challenge, however, is a structural mismatch between mission-level intelligence requirements and the resource limits of small UAVs. Although advanced AI models can improve perception and autonomy, small UAVs remain constrained in computation, energy, memory, and communication capability. This mismatch becomes more severe as modern deep learning models grow in size and complexity, creating a fundamental tension between the level of intelligence required for autonomous operation and the resources available at the aerial edge. Resolving this contradiction is the main motivation of this thesis \cite{hu2025splitreview,duan2022combined,li2024fsl,mcenroe2022survey,zhou2019edge,deng2020edge}.

Existing solution directions can be grouped into three broad categories. The first is onboard inference, in which models execute entirely on a single UAV and avoid network dependence, but remain tightly limited by SWaP constraints. The second is cloud or edge offloading, which shifts computation to more capable infrastructure and can reduce onboard burden, but depends on reliable connectivity and may incur unacceptable latency. The third is collaborative inference, in which model execution is distributed across multiple nodes; this direction is promising because it can exploit swarm resources collectively, yet it introduces tighter dependence between computation placement and communication cost \cite{qu2024elastic,wang2024jointuav,baktayan2024slr}.

These three directions expose a clear research gap. Pure onboard inference remains capability-limited, while cloud-based approaches are vulnerable to latency and connectivity disruption in remote or contested environments. Collaborative inference offers a stronger foundation, but its effectiveness depends on jointly optimizing model partitioning, node assignment, and communication cost, while adapting to mobility and failures. Existing studies rarely treat these aspects within one coherent optimization perspective. Two specific shortcomings stand out:

\begin{itemize}
  \item \textbf{The Coupling Gap.} Current approaches typically treat model partitioning and communication scheduling as loosely connected processes. In practice, however, partition boundaries directly determine the size of intermediate tensors exchanged between UAVs, which in turn govern transmission delay and RF energy consumption. This \emph{computation--communication coupling} means that partitioning and assignment decisions cannot be optimized independently without sacrificing global efficiency.

  \item \textbf{The Adaptation Gap.} Even when a jointly optimal configuration is found, it is only valid for the current system state. In a UAV swarm, link quality, residual energy, and topology change continuously due to mobility and potential node failures. A static solution therefore degrades over time. Existing online approaches re-decide partitioning only for newly arriving tasks; principled mechanisms that re-derive the coupled partitioning--assignment structure for inferences already in flight---and resume them without restarting---are still missing.
\end{itemize}

These two gaps motivate the research questions pursued in this thesis \cite{qu2024elastic,lansky2023firefly,hosseinzadeh2024local,pan2025selfhealing,jamshidi2025resilient}.
